\documentclass[11pt]{article}
\usepackage[a4paper,margin=27mm]{geometry}
\usepackage{amsmath,amssymb,mathtools}
\usepackage{fontspec,unicode-math}
\setmainfont{Latin Modern Roman}
\setmathfont{Latin Modern Math}
\usepackage{xurl,hyperref}
\hypersetup{colorlinks=true,urlcolor=blue,linkcolor=blue}
\newcommand{\sref}[2]{\hyperlink{source:#1:#2}{[S#2]}}
\newcommand{\cataloguescope}[1]{\paragraph{Mathematical question.} #1}
\begin{document}
% UnsolvedMath ID 8500010; source catalogue code AMR-084-0010
\setcounter{section}{8500009}
\section{Degree-only bounds for polynomial D(n)-tuples}\label{8500010}
{\small\sffamily UnsolvedMath ID 8500010 · Graph Theory}
\subsection{Definitions and mathematical statement}

\paragraph{Shared notation.}$\mathbb N=\{1,2,\ldots\}$, and $\mathbb Z,\mathbb Q,\mathbb R,\mathbb C$ denote the integers, rationals, reals and complex numbers. Any other conventions are specified locally.


For a nonzero $n\in\mathbb Z[X]$, call a set $A=\{a_1,\ldots,a_m\}$ of distinct nonzero polynomials in $\mathbb Z[X]$ a polynomial $D(n)$-tuple if $a_i a_j+n\in\{b(X)^2:b\in\mathbb Z[X]\}$ for all $i<j$, and there is no nonzero $p\in\mathbb Z[X]$ for which every $a_i/p$ and $n/p^2$ is a constant rational integer. This last condition excludes sets obtained by multiplying an integer tuple and its parameter by a common polynomial and its square. Let $P_n$ be the supremum of their cardinalities.
\paragraph{Statement recorded in the supplied source.}
Is there a function $B:\mathbb Z_{\ge0}\to\mathbb N$ such that $P_n\le B(\deg n)$ for every nonzero integer polynomial $n$, independent of the sizes of its coefficients? \sref{8500010}{2}
\subsection{Short English statement}
Can polynomial Diophantine tuples be bounded in size solely by the degree of the parameter polynomial?
\subsection{Sources}
\begin{description}
\item[\hypertarget{source:8500010:1}{S1}] ULAM.AI and the UnsolvedMath Contributors, UnsolvedMath: A Curated Collection of Open Mathematics Problems, record 8500010 (AMR-084-0010). CC BY 4.0. \url{https://huggingface.co/datasets/ulamai/UnsolvedMath}.
\item[\hypertarget{source:8500010:2}{S2}] Andrej Dujella, Diophantine m-tuples, author-maintained survey (snapshot accessed 9 October 2026), Section 7.2, definition of polynomial tuples and degree-only upper-bound question. \url{https://raw.githubusercontent.com/dujella/dujella.github.io/master/gener.html}.
\end{description}
\label{end:8500010}
\end{document}
